% !TeX program = pdflatex
% Converted from Markdown for Overleaf/pdfLaTeX.
% This file intentionally avoids fontspec and is compatible with pdfLaTeX.
\documentclass[12pt,a4paper]{article}

\usepackage[utf8]{inputenc}
\usepackage{cmap}
\usepackage[T2A]{fontenc}
\usepackage[english,bulgarian]{babel}
\selectlanguage{bulgarian}
\usepackage{geometry}
\geometry{margin=2.5cm}
\usepackage{microtype}
\usepackage{setspace}
\onehalfspacing
\setlength{\parindent}{0pt}
\setlength{\parskip}{6pt plus 2pt}
\usepackage{xcolor}
\usepackage{graphicx}
\usepackage{float}
\usepackage{caption}
\usepackage{amsmath}
\usepackage{longtable,booktabs,array,multirow}
\usepackage{pdflscape}
\usepackage{etoolbox}
\setlength{\LTleft}{0pt}
\setlength{\LTright}{0pt}
\setlength{\LTcapwidth}{\textwidth}
\usepackage{enumitem}
\setlist{itemsep=2pt,topsep=4pt}
\usepackage{fancyvrb}
\usepackage{fvextra}
\DefineVerbatimEnvironment{verbatim}{Verbatim}{breaklines=true,breakanywhere=true,fontsize=\small}

\providecommand{\tightlist}{\setlength{\itemsep}{0pt}\setlength{\parskip}{0pt}}
\providecommand{\pandocbounded}[1]{#1}
\providecommand{\passthrough}[1]{#1}
\sloppy
\hbadness=10000
\hfuzz=1000pt

\newcommand{\code}[1]{\texttt{#1}}
\newcommand{\metric}[1]{\textbf{#1}}
\newcommand{\todo}[1]{\textcolor{red}{#1}}

\begin{document}

\begin{titlepage}
\centering
\vspace*{1.2cm}
{\Large Софийски университет ``Св. Климент Охридски''\\[0.3cm]}
{\Large Факултет по математика и информатика\\[0.3cm]}
{\large Предмет: Въведение в дълбокото самообучение\\[1.5cm]}

{\Huge\bfseries Тема: 3. Семантично сегментиране на снимки\\[0.8cm]}
{\Large Проектна документация\\[1.5cm]}

\begin{tabular}{rl}
Изготвил: & Явор Чамов \\
Факултетен номер: & 4MI3400787 \\
Финален модел: & FPN + EfficientNet-B3 + CE Dice \\
\end{tabular}

\vfill
{\large \today}
\end{titlepage}

\tableofcontents
\newpage

\section{Резюме}

Настоящият документ описва цялостната реализация на проект за \textbf{семантично сегментиране} върху \textbf{Cityscapes Dataset}. Целта е да се обучат и сравнят различни невронни архитектури, които приемат RGB изображение от градска сцена и връщат маска, в която всеки пиксел е класифициран към един от 19 семантични класа.

Проектът следва пълен работен процес: подготовка на данни, exploratory data analysis, имплементация на training pipeline, сравнение на архитектури, loss функции и learning-rate schedulers, изграждане на model report и Streamlit demo приложение. Основният план за реализация включва изисквания за Cityscapes dataset, сравнение на Cross Entropy, Dice и Focal loss, експерименти с autoencoder/U-Net тип мрежа, cosine annealing и step decay scheduler-и, Streamlit интерфейс и tests.

Финално най-добър модел е избран \code{046\_fpn\_efficientnet\_b3\_ce\_dice}, който използва \textbf{FPN архитектура}, \textbf{EfficientNet-B3 encoder}, ImageNet pretrained weights, \textbf{Cross Entropy + Dice loss} и \textbf{Cosine Annealing scheduler}. Моделът постига:

\begin{itemize}
  \item \metric{mean IoU} = 0.7283;
  \item \metric{mean Dice} = 0.8335;
  \item \metric{pixel accuracy} = 0.9503.
\end{itemize}

\section{Въведение}

Семантичното сегментиране е задача от компютърното зрение, при която всеки пиксел от изображението се класифицира към предварително дефиниран клас. За разлика от image classification, където се предсказва един клас за цялото изображение, и object detection, където се предсказват bounding boxes, semantic segmentation връща плътна двумерна карта на класовете.

Тази задача е особено важна при автономно шофиране и анализ на градски сцени, защото системата трябва да разбира къде се намират пътят, тротоарът, автомобилите, пешеходците, знаците, светофарите и други обекти. В контекста на Cityscapes целта е да се класифицират пикселите към 19 тренировъчни класа.

\section{Цел на проекта и изисквания}

Основната цел е изграждане на end-to-end система за semantic segmentation на градски сцени. Системата трябва да позволява обучение, оценяване, сравнение и демонстрация на различни модели.

Проектът покрива следните изисквания:

\begin{enumerate}
  \item Проучване на научни статии и техники за semantic segmentation.
  \item Сравнение на pixel-wise Cross Entropy loss, Dice loss и Focal loss.
  \item Анализ на Cityscapes dataset: брой изображения, класове, class imbalance, аномалии и визуализации.
  \item Експерименти с autoencoder-type архитектура, реализирана чрез Tiny U-Net.
  \item Демонстрация на learning-rate scheduling чрез cosine annealing и step decay.
  \item Експерименти с различни preprocessing и augmentation стратегии.
  \item Streamlit web интерфейс за inference.
  \item Unit tests за ключови компоненти.
  \item Model report файл с baseline, експерименти, метрики, коментари и графики.
\end{enumerate}

\section{Dataset: Cityscapes}

\subsection{Използвани пакети}

За проекта са нужни два основни Cityscapes пакета:

\begin{verbatim}
leftImg8bit_trainvaltest.zip
където са RGB изображенията за train, val и test

gtFine_trainvaltest.zip
където са фините pixel-level анотации за train и val
\end{verbatim}

Очакваната структура на данните е:

\begin{verbatim}
data/raw/cityscapes/
+-- leftImg8bit/
|   +-- train/
|   +-- val/
|   +-- test/
+-- gtFine/
    +-- train/
    +-- val/
    +-- test/
\end{verbatim}

\subsection{Използване на split-овете}

\begin{itemize}
  \item \code{train} -- използва се за обучение на моделите;
  \item \code{val} -- използва се за сравнение на моделите и финална оценка;
  \item \code{test} -- използва се само за inference/demo, защото публичният test split няма реални ground-truth masks.
\end{itemize}

\subsection{Класове}

В Cityscapes се използват 19 тренировъчни класа. Те са:

\begin{longtable}{p{0.25\textwidth} p{0.25\textwidth} p{0.25\textwidth}}
\toprule
Пътна инфраструктура & Обекти / среда & Участници в движението \\
\midrule
road & building & person \\
sidewalk & wall & rider \\
terrain & fence & car \\
 & pole & truck \\
 & traffic light & bus \\
 & traffic sign & train \\
 & vegetation & motorcycle \\
 & sky & bicycle \\
\bottomrule
\end{longtable}

Част от оригиналните Cityscapes label IDs се игнорират при обучение. За тази цел се използва \code{labelTrainIds} mapping и \code{ignore\_index = 255} за невалидни или неизползвани пиксели.

\section{Exploratory Data Analysis и подготовка на данните}

EDA фазата има за цел да провери дали dataset-ът е консистентен и подходящ за обучение. Реализационният план изисква анализ на броя наблюдения, class distribution, missing image-mask pairs, invalid mask values и визуализации.

Основните проверки са:

\begin{itemize}
  \item проверка на броя изображения и маски в \code{train} и \code{val};
  \item проверка дали всяко изображение има съответна маска;
  \item проверка на размерите на image-mask двойките;
  \item проверка на валидните train IDs и \code{255} ignore пиксели;
  \item анализ на class imbalance;
  \item визуализация на image, ground truth mask и overlay.
\end{itemize}

\subsection{Preprocessing}

Използваната основна preprocessing конфигурация е:

\begin{itemize}
  \item resize на изображението до \code{512x1024};
  \item random crop до \code{512x512};
  \item ImageNet normalization;
  \item синхронизирани геометрични трансформации върху image и mask;
  \item nearest-neighbor interpolation за маските.
\end{itemize}

Важно е маските да не се нормализират като изображения. Те трябва да останат като цели числа, представящи класове.

\subsection{Augmentation стратегии}

Тествани са няколко preprocessing/augmentation стратегии:

\begin{itemize}
  \item \code{resize\_only} -- само преоразмеряване;
  \item \code{basic\_aug} -- resize/crop, horizontal flip, brightness/contrast;
  \item \code{strong\_aug} -- по-силни цветови и геометрични трансформации.
\end{itemize}

Изводът от експериментите е, че по-силните трансформации не винаги подобряват mean IoU. В някои случаи те влошават стабилността на validation резултатите.

\begin{figure}[H]
\centering
\includegraphics[width=0.92\textwidth]{reports/presentation_assets/class_distribution.png}
\caption{Разпределение на пикселите по класове. Графиката показва силен class imbalance, доминиран от road, building и vegetation.}
\end{figure}

\begin{figure}[H]
\centering
\includegraphics[width=0.82\textwidth]{reports/presentation_assets/category_distribution.png}
\caption{Разпределение на пикселите по семантични категории. Flat, construction и nature категориите покриват най-голяма част от сцените.}
\end{figure}

\subsection{EDA визуализации от \code{reports/figures}}

Следващите фигури са директно генерирани от \code{make eda} и са запазени в \code{reports/figures}. Те документират както количествените характеристики на dataset-а, така и качествени примери за image-mask съответствие.

\begin{figure}[H]
\centering
\includegraphics[width=0.9\textwidth]{reports/figures/image_size_distribution.png}
\caption{Разпределение на размерите на изображенията в анализираната извадка. Cityscapes изображенията са консистентни като резолюция.}
\end{figure}

\begin{figure}[H]
\centering
\includegraphics[width=0.92\textwidth]{reports/figures/class_distribution.png}
\caption{Class distribution, генериран от EDA скрипта. Графиката потвърждава силния дисбаланс в полза на големите структурни класове.}
\end{figure}

\begin{figure}[H]
\centering
\includegraphics[width=0.92\textwidth]{reports/figures/category_pixel_proportions.png}
\caption{Дял на пикселите по семантични категории. Flat и construction категориите доминират в градските сцени.}
\end{figure}

\begin{figure}[H]
\centering
\includegraphics[width=0.94\textwidth]{reports/figures/paper_class_pixels_by_category.png}
\caption{Брой пиксели по класове, групирани по категории. Тази визуализация следва стила на dataset анализите в semantic segmentation статии.}
\end{figure}

\begin{figure}[H]
\centering
\includegraphics[width=0.92\textwidth]{reports/figures/annotation_summary_table.png}
\caption{Обобщение на анотациите и label density. Таблицата служи като compact EDA преглед на dataset-а.}
\end{figure}

\begin{figure}[H]
\centering
\includegraphics[width=0.94\textwidth]{reports/figures/sample_overlays.png}
\caption{Примерни image-mask overlays. Фигурата показва, че маските са правилно подравнени с RGB изображенията.}
\end{figure}

\begin{figure}[H]
\centering
\includegraphics[width=0.94\textwidth]{reports/figures/rare_classes_examples.png}
\caption{Примери с редки класове. Тези изображения са важни за qualitative анализ, защото редките класове често са най-трудни за сегментиране.}
\end{figure}

\begin{figure}[H]
\centering
\includegraphics[width=0.88\textwidth]{reports/figures/traffic_participant_complexity.png}
\caption{Оценка на scene complexity чрез traffic participant компоненти. По-сложните сцени съдържат повече малки и припокриващи се обекти.}
\end{figure}

\begin{figure}[H]
\centering
\includegraphics[width=0.88\textwidth]{reports/figures/vehicle_component_area_histogram.png}
\caption{Хистограма на площите на vehicle компоненти. Малките компоненти обясняват защо car, truck, bus и motorcycle класовете изискват добра multi-scale обработка.}
\end{figure}

\section{Training pipeline}

Проектът използва конфигурируем training pipeline, управляван чрез YAML конфигурации. Това позволява повторяемост и ясно сравнение между експериментите.

Основните компоненти са:

\begin{enumerate}
  \item Dataset loader за Cityscapes.
  \item Preprocessing и augmentations.
  \item Model factory за избор на архитектура.
  \item Loss factory за избор на loss функция.
  \item Optimizer и scheduler.
  \item Training loop с mixed precision.
  \item Evaluation loop с mean IoU, Dice и pixel accuracy.
  \item Checkpoint saving.
  \item Export на history, plots, confusion matrix и model report.
\end{enumerate}

Примерна конфигурация:

\begin{verbatim}
training:
  epochs: 30
  batch_size: 4
  mixed_precision: true

model:
  architecture: deeplabv3plus
  encoder_name: resnet101
  encoder_weights: imagenet
  num_classes: 19

loss:
  name: cross_entropy_dice

scheduler:
  name: cosine_annealing
\end{verbatim}

\section{Използвани архитектури}

\subsection{Baseline majority-class model}

Baseline моделът е greedy статистически модел, който предсказва най-често срещания клас за всеки валиден пиксел. В проекта този модел е \code{000\_baseline\_majority}.

Неговите резултати са:

\begin{itemize}
  \item mean IoU = 0.0198;
  \item mean Dice = 0.0288;
  \item pixel accuracy = 0.3765.
\end{itemize}

Baseline е слаб, но е важен, защото всички невронни модели трябва да се сравняват спрямо него.

\subsection{Tiny U-Net}

Tiny U-Net е малка custom encoder-decoder архитектура, реализирана ръчно в PyTorch. Тя изпълнява изискването за autoencoder-type neural network. Архитектурата съдържа:

\begin{itemize}
  \item encoder blocks с convolution, batch normalization и ReLU;
  \item max pooling за downsampling;
  \item bottleneck слой;
  \item decoder с upsampling или transposed convolution;
  \item skip connections чрез concatenation;
  \item final \code{1x1} convolution към 19 класа.
\end{itemize}

Tiny U-Net е полезен за валидиране на pipeline-а, но значително изостава спрямо pretrained encoder моделите.

\subsection{U-Net}

U-Net е класическа encoder-decoder архитектура за segmentation. Encoder-ът извлича все по-абстрактни признаци, а decoder-ът възстановява пространствената резолюция. Skip connections прехвърлят детайлна локална информация от encoder към decoder.

В проекта U-Net е използван с различни encoders:

\begin{itemize}
  \item ResNet18;
  \item ResNet34;
  \item ResNet50;
  \item EfficientNet-B3.
\end{itemize}

Най-добрата U-Net конфигурация е \code{030\_final\_unet\_resnet34\_ce\_dice\_cosine}, която постига mIoU = 0.6920.

\subsection{U-Net++}

U-Net++ разширява U-Net чрез по-гъсти skip connections между encoder и decoder блоковете. Целта е да се намали semantic gap между ниско-level feature maps и decoder feature maps.

В проекта \code{044\_unetplusplus\_resnet34\_ce\_dice} постига mIoU = 0.6925, което е малко над стандартния U-Net ResNet34 CE+Dice резултат.

\subsection{FPN}

Feature Pyramid Network (FPN) използва признаци от различни дълбочини на encoder-а и ги агрегира в multi-scale представяне. Това е особено полезно за Cityscapes, защото в една сцена има големи обекти като road/building и малки обекти като traffic light/sign.

Финалният най-добър модел е \code{046\_fpn\_efficientnet\_b3\_ce\_dice} с mIoU = 0.7283.

\subsection{PSPNet}

PSPNet използва pyramid pooling за улавяне на глобален контекст. Идеята е моделът да разбира по-добре цялата сцена, а не само локалните детайли.

В проекта \code{024\_final\_pspnet\_resnet34\_ce\_cosine} постига mIoU = 0.5835. Това е по-ниско от FPN, DeepLabV3+ и U-Net конфигурациите, но е полезен сравнителен експеримент.

\subsection{DeepLabV3+}

DeepLabV3+ комбинира atrous convolution, Atrous Spatial Pyramid Pooling (ASPP) и decoder модул. Atrous convolution позволява по-голям receptive field без директно намаляване на резолюцията.

Най-добрият DeepLabV3+ експеримент е \code{047\_deeplabv3plus\_efficientnet\_b3\_ce\_dice} с mIoU = 0.7256. Той е много близък до най-добрия FPN модел.

\subsection{SegFormer}

SegFormer е transformer-based архитектура за semantic segmentation. Тя използва hierarchical transformer encoder и лек MLP decoder. В проекта \code{049\_segformer\_mit\_b1\_ce\_dice\_optional} е най-добрият transformer-based експеримент с mIoU = 0.7087.

\section{Използвани encoders}

\subsection{ResNet}

ResNet архитектурите използват residual connections, които улесняват обучението на по-дълбоки мрежи. В проекта са използвани ResNet18, ResNet34, ResNet50 и ResNet101.

Разликата между тях е основно в дълбочината и броя параметри. По-дълбоките модели могат да извличат по-богати признаци, но изискват повече изчислителни ресурси и не винаги дават пропорционално подобрение.

\subsection{EfficientNet-B3}

EfficientNet използва compound scaling, при който depth, width и resolution се увеличават балансирано. EfficientNet-B3 дава по-силни pretrained features от ResNet34 в най-добрите експерименти.

Най-добрите два резултата използват EfficientNet-B3 encoder:

\begin{itemize}
  \item FPN + EfficientNet-B3: mIoU = 0.7283;
  \item DeepLabV3+ + EfficientNet-B3: mIoU = 0.7256.
\end{itemize}

\subsection{MiT-B1}

MiT-B1 е encoder от SegFormer архитектурата. Той използва transformer blocks и hierarchical feature representation. В проекта MiT-B1 дава добър резултат, но не надминава EfficientNet-B3 моделите.

\section{Loss функции}

\subsection{Cross Entropy Loss}

Pixel-wise Cross Entropy разглежда всеки пиксел като отделен classification пример. За multi-class segmentation моделът връща logits с форма \code{[B, C, H, W]}, а target маската е \code{[B, H, W]}.

\[
L_{CE} = - \sum_{c=1}^{C} y_c \log(p_c)
\]

Предимството на Cross Entropy е стабилното обучение. Недостатъкът е, че при силен class imbalance тя може да фаворизира доминиращите класове.

\subsection{Dice Loss}

Dice Loss оптимизира припокриването между predicted mask и ground truth mask.

\[
Dice = \frac{2TP}{2TP + FP + FN}
\]

Тя е полезна при segmentation, защото измерва overlap quality. Самостоятелно Dice Loss може да бъде по-нестабилна, затова в проекта най-добри резултати дава комбинацията Cross Entropy + Dice.

\subsection{Focal Loss}

Focal Loss намалява тежестта на лесните примери и се фокусира върху трудни пиксели и редки класове.

\[
L_{focal} = - (1 - p_t)^\gamma \log(p_t)
\]

В проекта Focal+Dice дава силни резултати, но не надминава CE+Dice.

\subsection{Lovasz Loss}

Lovasz е loss, ориентиран към IoU/Jaccard оптимизация. В проекта CE+Lovasz е тестван чрез \code{048\_unet\_resnet34\_ce\_lovasz}, но резултатът не надминава CE+Dice.

\section{Метрики за оценка}

\subsection{Pixel Accuracy}

Pixel Accuracy измерва дела на правилно класифицираните пиксели:

\[
PixelAccuracy = \frac{correct\ pixels}{all\ valid\ pixels}
\]

Тази метрика е лесна за разбиране, но не е достатъчна сама по себе си, защото Cityscapes има доминиращи класове като road, building и sky. Модел може да има висока pixel accuracy, но слаб IoU върху редки класове.

\subsection{Intersection over Union}

IoU е основната метрика в проекта:

\[
IoU = \frac{TP}{TP + FP + FN}
\]

Mean IoU е средната стойност на IoU по класове. Тя е по-информативна от pixel accuracy при class imbalance.

\subsection{Dice coefficient}

Dice coefficient е overlap метрика, близка по идея до IoU:

\[
Dice = \frac{2TP}{2TP + FP + FN}
\]

В model report се използват максимум три основни метрики: mean IoU, mean Dice и pixel accuracy.

\section{Експериментален план}

Първоначалният implementation plan предвижда baseline, tiny U-Net, U-Net, DeepLabV3+, FPN, PSPNet, U-Net++ и опционален SegFormer. Всички експерименти се управляват чрез YAML configs, а резултатите се записват в model report.

В проекта са реализирани общо 28 evaluation експеримента плюс baseline. Това надвишава минималното изискване от 8 експеримента и позволява надеждно сравнение между архитектури, encoders, loss функции, scheduler-и и preprocessing варианти.

\section{Model report и резултати}

Model report файлът съдържа:

\begin{itemize}
  \item baseline row като първи ред;
  \item всички evaluation експерименти;
  \item Top 10 highlight;
  \item процентна промяна спрямо baseline;
  \item коментари за всеки модел;
  \item loss curves за Top 5;
  \item artifacts/checkpoints sheet.
\end{itemize}


\begin{landscape}
\begingroup
\scriptsize
\setlength{\tabcolsep}{3pt}
\begin{longtable}{r p{7.0cm} p{3.0cm} p{3.2cm} p{3.2cm} r r r}
\caption{Най-добри 10 evaluation експеримента, сортирани по mean IoU.}\\
\toprule
Rank & Experiment & Architecture & Encoder & Loss & mIoU & Dice & Acc. \\
\midrule
\endfirsthead
\toprule
Rank & Experiment & Architecture & Encoder & Loss & mIoU & Dice & Acc. \\
\midrule
\endhead
1 & \texttt{046\_fpn\_efficientnet\_b3\_ce\_dice} & fpn & efficientnet-b3 & cross\_entropy\_dice & 0.7283 & 0.8335 & 0.9503 \\
2 & \texttt{047\_deeplabv3plus\_efficientnet\_b3\_ce\_dice} & deeplabv3plus & efficientnet-b3 & cross\_entropy\_dice & 0.7256 & 0.8310 & 0.9482 \\
3 & \texttt{049\_segformer\_mit\_b1\_ce\_dice\_optional} & segformer & mit\_b1 & cross\_entropy\_dice & 0.7087 & 0.8188 & 0.9488 \\
4 & \texttt{043\_deeplabv3plus\_resnet101\_ce\_dice} & deeplabv3plus & resnet101 & cross\_entropy\_dice & 0.7022 & 0.8130 & 0.9450 \\
5 & \texttt{044\_unetplusplus\_resnet34\_ce\_dice} & unetplusplus & resnet34 & cross\_entropy\_dice & 0.6925 & 0.8057 & 0.9495 \\
6 & \texttt{030\_final\_unet\_resnet34\_ce\_dice\_cosine} & unet & resnet34 & cross\_entropy\_dice & 0.6920 & 0.8066 & 0.9466 \\
7 & \texttt{040\_unet\_resnet34\_ce\_dice\_50ep} & unet & resnet34 & cross\_entropy\_dice & 0.6910 & 0.8045 & 0.9478 \\
8 & \texttt{045\_unet\_efficientnet\_b3\_ce\_dice} & unet & efficientnet-b3 & cross\_entropy\_dice & 0.6908 & 0.7946 & 0.9510 \\
9 & \texttt{031\_final\_unet\_resnet34\_focal\_dice\_cosine} & unet & resnet34 & focal\_dice & 0.6869 & 0.8020 & 0.9456 \\
10 & \texttt{048\_unet\_resnet34\_ce\_lovasz} & unet & resnet34 & cross\_entropy\_lovasz & 0.6865 & 0.8006 & 0.9476 \\

\bottomrule
\end{longtable}
\endgroup
\end{landscape}

\begin{figure}[H]
\centering
\includegraphics[width=0.92\textwidth]{reports/presentation_assets/top_models_miou.png}
\caption{Top evaluation runs, сортирани по validation mean IoU. Най-добрият модел е FPN + EfficientNet-B3 + CE+Dice.}
\end{figure}

\section{Избор на най-добър модел}

Най-добрият модел е \code{046\_fpn\_efficientnet\_b3\_ce\_dice}.

\begin{longtable}{p{0.35\textwidth} p{0.55\textwidth}}
\toprule
Компонент & Стойност \\
\midrule
Architecture & FPN \\
Encoder & EfficientNet-B3 \\
Encoder weights & ImageNet \\
Loss & Cross Entropy + Dice \\
Scheduler & Cosine Annealing \\
Epochs & 30 \\
Batch size & 4 \\
Image size / crop & 512x1024 / 512x512 \\
Augmentation & basic augmentation \\
Mean IoU & 0.7283 \\
Mean Dice & 0.8335 \\
Pixel Accuracy & 0.9503 \\
\bottomrule
\end{longtable}

Този модел работи най-добре поради комбинацията от:

\begin{itemize}
  \item multi-scale feature aggregation чрез FPN;
  \item силен pretrained EfficientNet-B3 encoder;
  \item CE+Dice loss, който балансира pixel-wise classification и overlap quality;
  \item Cosine Annealing scheduler, който стабилизира обучението.
\end{itemize}

\section{Анализ на top моделите}

\subsection{FPN + EfficientNet-B3}

FPN комбинира feature maps от различни нива на encoder-а. Това помага едновременно за големи класове като road и building и за по-малки обекти. EfficientNet-B3 предоставя по-силни представяния от ResNet34, което се вижда от най-високия mean IoU.

\subsection{DeepLabV3+ + EfficientNet-B3}

DeepLabV3+ с EfficientNet-B3 е вторият най-добър модел. Разликата спрямо FPN е малка, което показва, че atrous convolution и ASPP са много подходящи за Cityscapes.

\subsection{SegFormer MiT-B1}

SegFormer е най-добрият transformer-based run. Той се представя добре, но изостава от EfficientNet-B3 CNN моделите. Това може да се дължи на ограничен брой epochs, по-малък batch size или необходимост от по-фина настройка на learning rate.

\subsection{U-Net и U-Net++}

U-Net ResNet34 CE+Dice остава стабилен силен baseline. U-Net++ постига малко по-добър mIoU от стандартния U-Net, но не достига FPN и DeepLabV3+ с EfficientNet-B3.

\section{Loss curves и generalization gap}

За Top 5 моделите са анализирани train/validation loss curves. Основните наблюдения са:

\begin{itemize}
  \item train loss намалява стабилно при всички Top 5 модела;
  \item validation loss също намалява, но остава по-висок от train loss;
  \item има умерен generalization gap, но не драматично overfitting поведение;
  \item при някои модели validation loss е по-шумен;
  \item Top 1 и Top 2 са много близки като качество.
\end{itemize}

\section{Confusion matrix и error analysis}

За най-добрия модел е изготвена confusion matrix визуализация. Общите изводи от анализа са:

\begin{itemize}
  \item моделът е най-силен върху доминиращи и визуално стабилни класове като road, building, vegetation, car, sky и sidewalk;
  \item проблемните класове са малки, редки или boundary-sensitive;
  \item типични трудни класове са traffic light, traffic sign, rider, motorcycle и bicycle;
  \item чести грешки се появяват около граници между road, sidewalk и terrain;
  \item далечни и частично скрити обекти са по-трудни за сегментиране.
\end{itemize}

Това потвърждава защо mean IoU е по-подходяща основна метрика от pixel accuracy.

\begin{figure}[H]
\centering
\includegraphics[width=0.94\textwidth]{reports/presentation_assets/per_class_iou_bar.png}
\caption{Per-class IoU за най-добрия модел. Малките и редки класове имат по-нисък IoU от големите структурни класове.}
\end{figure}

\begin{figure}[H]
\centering
\includegraphics[width=0.94\textwidth]{reports/presentation_assets/per_class_recall_bar.png}
\caption{Per-class recall за най-добрия модел. Recall анализът показва кои класове моделът пропуска най-често.}
\end{figure}

\begin{figure}[H]
\centering
\includegraphics[width=0.94\textwidth]{reports/presentation_assets/confusion_matrix_normalized_by_true_class.png}
\caption{Нормализирана confusion matrix по истински клас. Диагоналът показва правилните класификации, а off-diagonal стойностите показват най-честите обърквания.}
\end{figure}

\begin{figure}[H]
\centering
\includegraphics[width=0.94\textwidth]{reports/presentation_assets/confusion_matrix_log_counts.png}
\caption{Confusion matrix с логаритмична скала на броя пиксели. Логаритмичната визуализация прави видими и по-редките грешки.}
\end{figure}

\begin{figure}[H]
\centering
\includegraphics[width=0.9\textwidth]{reports/presentation_assets/top_15_confusions_bar.png}
\caption{Top 15 най-чести обърквания между класове. Тези грешки са основа за qualitative error analysis и бъдещи подобрения.}
\end{figure}

\subsection{Qualitative validation examples}

Следващите панели са генерирани от evaluation pipeline-а за най-добрия модел \code{046\_fpn\_efficientnet\_b3\_ce\_dice}. Всеки панел съдържа original image, ground truth, prediction, overlay и error map.

\begin{figure}[H]
\centering
\includegraphics[width=0.96\textwidth]{reports/figures/046_fpn_efficientnet_b3_ce_dice_eval/example_000_panel.png}
\caption{Qualitative пример 1 за най-добрия модел. Панелът показва едновременно входното изображение, истинската маска, предсказанието, overlay и error map.}
\end{figure}

\begin{figure}[H]
\centering
\includegraphics[width=0.96\textwidth]{reports/figures/046_fpn_efficientnet_b3_ce_dice_eval/example_001_panel.png}
\caption{Qualitative пример 2 за най-добрия модел. Error map частта позволява бърза визуална проверка на проблемните пиксели.}
\end{figure}

\begin{figure}[H]
\centering
\includegraphics[width=0.96\textwidth]{reports/figures/046_fpn_efficientnet_b3_ce_dice_eval/example_002_panel.png}
\caption{Qualitative пример 3 за най-добрия модел. Такива примери допълват числените метрики с визуална интерпретация на грешките.}
\end{figure}

\section{Streamlit demo приложение}

Крайният проект включва Streamlit приложение за демонстрация на inference. Основните функционалности са:

\begin{itemize}
  \item upload на изображение;
  \item избор на checkpoint;
  \item inference с избрания модел;
  \item показване на original image, predicted mask и overlay;
  \item легенда с класовете;
  \item опция за сравнение с ground truth, ако е наличен.
\end{itemize}

Streamlit приложението демонстрира практическата използваемост на проекта и затваря цикъла от обучение до визуална демонстрация.

\section{Тестове и качество на кода}

Implementation plan предвижда unit tests, които не изискват пълния Cityscapes dataset и могат да се изпълняват на CPU. Основните тестове са:

\begin{itemize}
  \item label mapping тестове;
  \item dataset loading тестове със synthetic data;
  \item loss function тестове;
  \item metric тестове;
  \item model factory тестове;
  \item training step smoke test;
  \item inference и Streamlit import smoke tests.
\end{itemize}

Тези тестове гарантират, че pipeline-ът е възпроизводим и стабилен.

\section{Ограничения}

Проектът има няколко ограничения:

\begin{itemize}
  \item тренировката е извършена върху resize/crop изображения, а не върху full-resolution Cityscapes изображения;
  \item някои архитектури са ограничени от GPU memory;
  \item публичният test split няма реални ground-truth masks;
  \item class imbalance остава предизвикателство;
  \item SegFormer експериментите са ограничени до MiT-B1;
  \item качественият анализ зависи от избрани validation примери.
\end{itemize}

\section{Бъдеща работа}

Възможни подобрения:

\begin{itemize}
  \item class-weighted Cross Entropy;
  \item по-големи crops: 512x768 или 768x768;
  \item test-time augmentation;
  \item fine-tuning на най-добрия модел с по-нисък learning rate;
  \item SegFormer MiT-B2 или MiT-B3;
  \item DeepLabV3+ с други EfficientNet encoders;
  \item per-class confusion matrix и error analysis по класове;
  \item повече qualitative examples за добри и проблемни предсказания.
\end{itemize}

\section{Заключение}

Проектът демонстрира цялостен deep learning workflow за semantic segmentation: данни, preprocessing, модели, обучение, оценяване, model report и Streamlit demo.

Основният извод е, че \textbf{multi-scale feature aggregation}, комбиниран със \textbf{силен pretrained encoder} и \textbf{CE+Dice loss}, дава най-добро качество върху Cityscapes validation split. Най-добрият модел е FPN + EfficientNet-B3, който постига mean IoU = 0.7283.

Pixel Accuracy остава полезна допълнителна метрика, но не трябва да се използва самостоятелно, защото може да прикрие слабо представяне върху редки класове. Mean IoU е основната метрика за вземане на решение за най-добър модел.

\newpage
\appendix
\section{Пълен списък на експериментите}


\begin{landscape}
\begingroup
\scriptsize
\setlength{\tabcolsep}{3pt}
\begin{longtable}{r p{7.0cm} p{4.2cm} p{3.0cm} p{3.2cm} r r r}
\caption{Пълен списък на baseline и всички evaluation експерименти от model report файла.}\\
\toprule
Rank & Experiment ID & Model / Encoder & Loss & Scheduler & mIoU & Dice & Acc. \\
\midrule
\endfirsthead
\toprule
Rank & Experiment ID & Model / Encoder & Loss & Scheduler & mIoU & Dice & Acc. \\
\midrule
\endhead
0 & \texttt{000\_baseline\_majority} & majority\_baseline / - & cross\_entropy & cosine\_annealing & 0.0198 & 0.0288 & 0.3765 \\
1 & \texttt{046\_fpn\_efficientnet\_b3\_ce\_dice} & fpn / efficientnet-b3 & cross\_entropy\_dice & cosine\_annealing & 0.7283 & 0.8335 & 0.9503 \\
2 & \texttt{047\_deeplabv3plus\_efficientnet\_b3\_ce\_dice} & deeplabv3plus / efficientnet-b3 & cross\_entropy\_dice & cosine\_annealing & 0.7256 & 0.8310 & 0.9482 \\
3 & \texttt{049\_segformer\_mit\_b1\_ce\_dice\_optional} & segformer / mit\_b1 & cross\_entropy\_dice & cosine\_annealing & 0.7087 & 0.8188 & 0.9488 \\
4 & \texttt{043\_deeplabv3plus\_resnet101\_ce\_dice} & deeplabv3plus / resnet101 & cross\_entropy\_dice & cosine\_annealing & 0.7022 & 0.8130 & 0.9450 \\
5 & \texttt{044\_unetplusplus\_resnet34\_ce\_dice} & unetplusplus / resnet34 & cross\_entropy\_dice & cosine\_annealing & 0.6925 & 0.8057 & 0.9495 \\
6 & \texttt{030\_final\_unet\_resnet34\_ce\_dice\_cosine} & unet / resnet34 & cross\_entropy\_dice & cosine\_annealing & 0.6920 & 0.8066 & 0.9466 \\
7 & \texttt{040\_unet\_resnet34\_ce\_dice\_50ep} & unet / resnet34 & cross\_entropy\_dice & cosine\_annealing & 0.6910 & 0.8045 & 0.9478 \\
8 & \texttt{045\_unet\_efficientnet\_b3\_ce\_dice} & unet / efficientnet-b3 & cross\_entropy\_dice & cosine\_annealing & 0.6908 & 0.7946 & 0.9510 \\
9 & \texttt{031\_final\_unet\_resnet34\_focal\_dice\_cosine} & unet / resnet34 & focal\_dice & cosine\_annealing & 0.6869 & 0.8020 & 0.9456 \\
10 & \texttt{048\_unet\_resnet34\_ce\_lovasz} & unet / resnet34 & cross\_entropy\_lovasz & cosine\_annealing & 0.6865 & 0.8006 & 0.9476 \\
11 & \texttt{023\_final\_fpn\_resnet34\_ce\_cosine} & fpn / resnet34 & cross\_entropy & cosine\_annealing & 0.6838 & 0.8000 & 0.9468 \\
12 & \texttt{039\_final\_deeplabv3plus\_resnet101\_ce\_cosine} & deeplabv3plus / resnet101 & cross\_entropy & cosine\_annealing & 0.6793 & 0.7943 & 0.9480 \\
13 & \texttt{025\_final\_deeplabv3plus\_resnet50\_ce\_cosine} & deeplabv3plus / resnet50 & cross\_entropy & cosine\_annealing & 0.6730 & 0.7887 & 0.9472 \\
14 & \texttt{032\_final\_unet\_resnet34\_resize\_only\_ce\_cosine} & unet / resnet34 & cross\_entropy & cosine\_annealing & 0.6696 & 0.7871 & 0.9471 \\
15 & \texttt{022\_final\_unetplusplus\_resnet34\_ce\_cosine} & unetplusplus / resnet34 & cross\_entropy & cosine\_annealing & 0.6686 & 0.7857 & 0.9486 \\
16 & \texttt{033\_final\_unet\_resnet34\_strong\_aug\_ce\_cosine} & unet / resnet34 & cross\_entropy & cosine\_annealing & 0.6651 & 0.7835 & 0.9463 \\
17 & \texttt{021\_final\_unet\_resnet34\_ce\_cosine} & unet / resnet34 & cross\_entropy & cosine\_annealing & 0.6648 & 0.7828 & 0.9469 \\
18 & \texttt{038\_final\_unet\_resnet50\_ce\_cosine} & unet / resnet50 & cross\_entropy & cosine\_annealing & 0.6586 & 0.7763 & 0.9468 \\
19 & \texttt{027\_final\_unet\_resnet34\_ce\_plateau} & unet / resnet34 & cross\_entropy & reduce\_on\_plateau & 0.6558 & 0.7760 & 0.9446 \\
20 & \texttt{029\_final\_unet\_resnet34\_focal\_cosine} & unet / resnet34 & focal & cosine\_annealing & 0.6556 & 0.7752 & 0.9448 \\
21 & \texttt{037\_final\_unet\_resnet18\_ce\_cosine} & unet / resnet18 & cross\_entropy & cosine\_annealing & 0.6398 & 0.7603 & 0.9436 \\
22 & \texttt{026\_final\_unet\_resnet34\_ce\_step} & unet / resnet34 & cross\_entropy & step\_decay & 0.6037 & 0.7232 & 0.9419 \\
23 & \texttt{024\_final\_pspnet\_resnet34\_ce\_cosine} & pspnet / resnet34 & cross\_entropy & cosine\_annealing & 0.5835 & 0.7101 & 0.9305 \\
24 & \texttt{028\_final\_unet\_resnet34\_dice\_cosine} & unet / resnet34 & dice & cosine\_annealing & 0.5488 & 0.6363 & 0.9308 \\
25 & \texttt{036\_final\_tiny\_unet\_64\_ce\_cosine} & tiny\_unet / - & cross\_entropy & cosine\_annealing & 0.3920 & 0.4841 & 0.8808 \\
26 & \texttt{035\_final\_tiny\_unet\_32\_ce\_cosine} & tiny\_unet / - & cross\_entropy & cosine\_annealing & 0.3437 & 0.4294 & 0.8657 \\
27 & \texttt{020\_final\_tiny\_unet\_ce\_cosine} & tiny\_unet / - & cross\_entropy & cosine\_annealing & 0.3419 & 0.4265 & 0.8657 \\
28 & \texttt{034\_final\_tiny\_unet\_16\_ce\_cosine} & tiny\_unet / - & cross\_entropy & cosine\_annealing & 0.2610 & 0.3236 & 0.8323 \\

\bottomrule
\end{longtable}
\endgroup
\end{landscape}


\section{Примерна структура на проекта}

\begin{verbatim}
semantic-segmentation-cityscapes/
+-- app/
|   +-- streamlit_app.py
+-- configs/
|   +-- experiments/
+-- data/
|   +-- raw/cityscapes/
+-- docs/
+-- reports/
+-- scripts/
+-- src/cityseg/
|   +-- data/
|   +-- models/
|   +-- training/
|   +-- inference/
|   +-- reporting/
+-- tests/
\end{verbatim}

\section{Използвана литература}

\begin{enumerate}
  \item Cordts et al., The Cityscapes Dataset for Semantic Urban Scene Understanding, CVPR 2016. URL: \texttt{https://arxiv.org/abs/1604.01685}
  \item Long, Shelhamer, Darrell, Fully Convolutional Networks for Semantic Segmentation, CVPR 2015. URL: \texttt{https://arxiv.org/abs/1411.4038}
  \item Ronneberger, Fischer, Brox, U-Net: Convolutional Networks for Biomedical Image Segmentation, MICCAI 2015. URL: \texttt{https://arxiv.org/abs/1505.04597}
  \item Noh, Hong, Han, Learning Deconvolution Network for Semantic Segmentation, ICCV 2015. URL: \texttt{https://arxiv.org/abs/1505.04366}
  \item Chen et al., Encoder-Decoder with Atrous Separable Convolution for Semantic Image Segmentation, ECCV 2018. URL: \texttt{https://arxiv.org/abs/1802.02611}
  \item Zhao et al., Pyramid Scene Parsing Network, CVPR 2017. URL: \texttt{https://arxiv.org/abs/1612.01105}
  \item Xie et al., SegFormer: Simple and Efficient Design for Semantic Segmentation with Transformers, NeurIPS 2021. URL: \texttt{https://arxiv.org/abs/2105.15203}
  \item Lin et al., Focal Loss for Dense Object Detection, ICCV 2017. URL: \texttt{https://arxiv.org/abs/1708.02002}
  \item Milletari, Navab, Ahmadi, V-Net: Fully Convolutional Neural Networks for Volumetric Medical Image Segmentation, 2016. URL: \texttt{https://arxiv.org/abs/1606.04797}
  \item Sudre et al., Generalised Dice Overlap as a Deep Learning Loss Function for Highly Unbalanced Segmentations, 2017. URL: \texttt{https://arxiv.org/abs/1707.03237}
\end{enumerate}

\end{document}
